\documentclass[11pt]{article}
\usepackage[letterpaper,margin=1in]{geometry}
\usepackage[T1]{fontenc}
\usepackage[utf8]{inputenc}
\usepackage{lmodern,microtype}
\usepackage{amsmath,amssymb,mathtools,bm,slashed}
\usepackage{xcolor,enumitem,booktabs,array,needspace}
\usepackage[most]{tcolorbox}
\usepackage{fancyhdr,titlesec,parskip,hyperref}
\definecolor{qftblue}{HTML}{245A8D}
\definecolor{qftgreen}{HTML}{2F6B4F}
\definecolor{qftgray}{HTML}{F4F6F8}
\hypersetup{colorlinks=true,linkcolor=qftblue,urlcolor=qftblue,pdfauthor={Nima Lashkari}}
\pagestyle{fancy}\fancyhf{}
\lhead{PHYS 662: Quantum Field Theory I}\cfoot{\thepage}
\setlength{\headheight}{14pt}
\titleformat{\section}{\Large\bfseries\color{qftblue}}{\thesection.}{0.6em}{}
\titleformat{\subsection}{\large\bfseries\color{qftgreen}}{\thesubsection}{0.6em}{}
\setlist[enumerate]{leftmargin=2em,itemsep=0.25em,topsep=0.35em}
\setlist[itemize]{leftmargin=2em,itemsep=0.2em,topsep=0.35em}
\newcommand{\dd}{\mathrm d}
\newcommand{\ii}{\mathrm i}
\newcommand{\eps}{\varepsilon}
\newcommand{\tr}{\operatorname{tr}}
\newcommand{\Tr}{\operatorname{Tr}}
\newcommand{\diag}{\operatorname{diag}}
\newcommand{\Lag}{\mathcal L}
\newcommand{\D}{\mathcal D}
\newcommand{\vac}{\lvert0\rangle}
\newcommand{\vbra}{\langle0\rvert}
\newcommand{\Tord}{\mathrm T}
\newcommand{\ket}[1]{\lvert#1\rangle}
\newcommand{\bra}[1]{\langle#1\rvert}
\newtcolorbox{keyidea}[1][]{enhanced,breakable,colback=qftgray,colframe=qftblue,
 boxrule=0.7pt,arc=1.5mm,title={#1},fonttitle=\bfseries,coltitle=white}
\newtcolorbox{remarkbox}[1][]{enhanced,breakable,colback=white,colframe=qftgreen,
 boxrule=0.6pt,arc=1.5mm,title={#1},fonttitle=\bfseries,coltitle=white}
\newcommand{\intp}{\int\!\frac{\dd^4p}{(2\pi)^4}}
\newcommand{\MSbar}{\overline{\mathrm{MS}}}
\newcommand{\module}[1]{\clearpage\section{#1}}
\newcommand{\braket}[2]{\langle#1\mid#2\rangle}
\newcommand{\dPi}{\dd\Pi}
\newcommand{\cM}{\mathcal M}
\newcommand{\cH}{\mathcal H}
\newcommand{\cS}{\mathcal S}
\newcommand{\Zp}{Z_{\rm pole}}
\newcommand{\Kallen}{\mathcal K}
\newcommand{\lecturefront}[4]{%
\rhead{Lecture #1}\hypersetup{pdftitle={PHYS 662 Lecture #1: #2}}
\begin{document}\pagenumbering{gobble}
\begin{titlepage}\centering\vspace*{1.05in}
{\scshape\Large PHYS 662 --- Quantum Field Theory I\par}
\vspace{0.35in}{\Huge\bfseries\color{qftblue}Lecture #1\par}
\vspace{0.2in}{\LARGE\bfseries #2\par}
\vspace{0.2in}{\large\color{qftgreen}#3\par}
\vspace{0.6in}{\large Nima Lashkari\par}{\large Fall 2026\par}
\vfill\begin{minipage}{0.86\textwidth}\small #4\par\medskip
Course page: \url{https://nimalashkari.com/quantum-field-theory-2026-part-i/}
\end{minipage}\vspace{0.5in}\end{titlepage}
\pagenumbering{roman}\setcounter{tocdepth}{1}\tableofcontents
\vfill\begin{remarkbox}[Place in the course]
Lectures 21--22 developed interacting correlators and scalar loops.
Lecture 23 defines the scattering states and the S-matrix.
Lecture 24 converts amplitudes into rates and cross sections.
Lecture 25 derives the reduction formula connecting these subjects.
Homework 8 covers this three-lecture sequence.
\end{remarkbox}
\small\textbf{Reading and conventions.} These notes use the course slides and
D.~Harlow, \href{https://www.mit.edu/~harlow/HarlowQFT1.pdf}{\textit{QFT I}},
Sections 8--9, with our covariant state normalization and mostly-minus metric.
Useful online companions are D.~Tong,
\href{https://www.damtp.cam.ac.uk/user/tong/qft/qfthtml/S3.html}{\textit{QFT}, Sections 3.6--3.7},
and D.~Buchholz and W.~Dybalski,
\href{https://arxiv.org/abs/math-ph/0509047}{\textit{Scattering in relativistic quantum field theory}},
Sections 1--3, for the scope of the asymptotic construction.
For textbook study, see Weinberg, \textit{The Quantum Theory of Fields},
Vol.~I, Chapter 3, and Peskin--Schroeder, \textit{An Introduction to QFT},
Chapters 4 and 7. The accompanying editorial memo records the online
sources consulted and the corrections to the slides.
\normalsize\clearpage\pagenumbering{arabic}}
\lecturefront{23}{Multiparticle States and the Scattering Matrix}
{What comes in, what goes out, and what a particle means}
{A scattering amplitude compares two descriptions of the same interacting
system. We first identify stable particles in the exact spectrum, construct
asymptotic multiparticle states, and then define the S-matrix. This order
separates physical assumptions from perturbative calculational tools.}

\section{Stable particles are properties of the exact theory}
A field is a local operator; a particle is a particular kind of state.
The two are related, but they are not interchangeable. Let \(\ket{\Omega}\)
be the unique translation-invariant vacuum and choose its energy to be zero.
A stable scalar particle of mass \(M>0\) has exact states satisfying
\begin{equation}
 P^\mu\ket{\bm p}=p^\mu\ket{\bm p},\qquad
 p^0=E_{\bm p}=\sqrt{\bm p^2+M^2},\qquad
 \braket{\bm p}{\bm q}=2E_{\bm p}(2\pi)^3\delta^{(3)}(\bm p-\bm q).
 \label{eq:oneparticle}
\end{equation}
Here and below \(\braket{a}{b}\) denotes an inner product.
The mass is the physical mass, not necessarily the parameter \(m(\mu)\)
in a renormalized Lagrangian. For spin, add a spin label and a Kronecker delta.
Our conventions remain \(\eta=\diag(1,-1,-1,-1)\), \(\hbar=c=1\).

A local scalar with the correct quantum numbers can have a nonzero overlap
\begin{equation}
 \bra{\Omega}\phi(0)\ket{\bm p}=\sqrt{\Zp},\qquad
 \bra{\Omega}\phi(x)\ket{\bm p}=\sqrt{\Zp}\,e^{-\ii p\cdot x}.
 \label{eq:overlap}
\end{equation}
We choose its phase so that this overlap is real and positive, and subtract
\(\langle\phi\rangle\) if necessary. Lorentz invariance makes the scalar
coefficient independent of \(\bm p\). An arbitrary local field may also
create multiparticle states. Isolating a particle requires a spectral
projection onto its mass shell, or the asymptotic procedure developed below.

Even the free-field expression \(\phi(x)\phi(y)\vac\) is not a pure
two-particle state. Wick's identity gives
\begin{equation}
 \phi(x)\phi(y)\vac=
 \langle0|\phi(x)\phi(y)|0\rangle\vac
 +:\!\phi(x)\phi(y)\!:\vac .
\end{equation}
The first term is a vacuum component. In an interacting theory normal
ordering by itself does not project onto an exact particle-number sector.

\begin{keyidea}[The assumption behind ordinary particle scattering]
We begin with stable massive particles and isolated mass shells, in a
local theory with sufficiently short-range interactions. These assumptions
permit well-separated incoming and outgoing packets. We do not assume
that every local field creates just one particle, or that particle number
is conserved during the collision.
\end{keyidea}
Stable bound states can also be particle species. Their internal structure
changes their form factors and interactions, not their status as possible
external particles. An unstable resonance does not give an exact stable
one-particle state; its use in decay calculations needs a separate
approximation, discussed in Lecture 24.

\module{The asymptotic Fock space and its normalization}
Build an auxiliary space \(\cH_{\rm as}\) from all stable species relevant
to the scattering experiment. For a scalar species use
\begin{equation}
 \dPi_p=\frac{\dd^3p}{(2\pi)^3\,2E_{\bm p}},\qquad
 [a(\bm p),a^\dagger(\bm q)]=2E_{\bm p}(2\pi)^3\delta^{(3)}(\bm p-\bm q).
\end{equation}
The Hamiltonian \(H_{\rm as}\) adds the \emph{physical} single-particle
energies. It describes independent asymptotic particles, not the bare
quanta of the quadratic part of an arbitrary Lagrangian split.

For identical bosons define
\(\ket{p_1,\ldots,p_n}=a^\dagger(p_1)\cdots a^\dagger(p_n)\vac_{\rm as}\).
Their inner product is
\begin{equation}
 \langle q_1,\ldots,q_n|p_1,\ldots,p_n\rangle
 =\sum_{\pi\in S_n}\prod_{j=1}^n
 2E_{p_j}(2\pi)^3\delta^{(3)}(\bm q_{\pi(j)}-\bm p_j).
 \label{eq:multinorm}
\end{equation}
For identical fermions each permutation carries its parity sign.
Integration over all ordered momenta therefore counts the same physical
bosonic or fermionic configuration \(n!\) times. The identity is
\begin{equation}
 1_{\rm as}=\sum_{\{n_a\}}\frac{1}{\prod_a n_a!}
 \sum_{\{s_j\}}\int\prod_j\dPi_{p_j}\,
 \ket{\{p_j,s_j,a_j\}}\bra{\{p_j,s_j,a_j\}}.
 \label{eq:complete}
\end{equation}
The sum includes the vacuum. There is no factorial between distinguishable
species. Write \(\cS_f=\prod_a n_a!\) for the final-state counting factor;
this is distinct from the operator \(S\).

Plane waves are distributional basis vectors. A normalized one-particle
packet is instead
\begin{equation}
 \ket{f}=\int\dPi_p\,f(\bm p)\ket{\bm p},\qquad
 \int\dPi_p\,|f(\bm p)|^2=1.
\end{equation}
A two-boson packet made with normalized functions \(f,g\) has squared norm
\(1+|(f,g)|^2\), where \((f,g)=\int\dPi_p f^*g\). When their momentum
supports are disjoint this norm is one. Exchange symmetry is already in
the state; we do not add a second exchange amplitude after symmetrizing.

\begin{remarkbox}[Why the factorial matters]
Act with the \(n=2\) term of \eqref{eq:complete} on \(\ket{p,q}\).
The direct and exchanged delta functions each reproduce \(\ket{p,q}\).
The factor \(1/2!\) makes their sum the identity. The same factor will
appear in a rate integrated over the full labeled two-particle phase space.
\end{remarkbox}

\module{Spectral decomposition of a two-point function}
The exact two-point function shows whether a chosen field detects a
particle. Take a Hermitian scalar with \(\langle\phi\rangle=0\).
Insert a complete set of exact energy-momentum states into the Wightman
function (there is no time ordering yet):
\begin{equation}
 W_2(x)=\bra{\Omega}\phi(x)\phi(0)\ket{\Omega}
 =\sum_n\int\dd\mu(n)\,e^{-\ii p_n\cdot x}
 |\bra{\Omega}\phi(0)\ket{n}|^2.
 \label{eq:wightman}
\end{equation}
The measure \(\dd\mu(n)\) includes continuum integrations and degeneracies.
This spectral insertion does not require a prior assumption of asymptotic
completeness. Translation invariance gave the exponential; positivity of
the Hilbert-space inner product gives nonnegative weights.

Define the scalar spectral measure by
\begin{equation}
 \sum_n\int\dd\mu(n)\,(2\pi)^4\delta^{(4)}(q-p_n)
 |\bra{\Omega}\phi(0)\ket{n}|^2
 =2\pi\theta(q^0)\rho(q^2).
 \label{eq:rhodef}
\end{equation}
The spectrum condition and Lorentz invariance then imply
\begin{align}
 W_2(x)&=\int_0^\infty\dd\mu^2\,\rho(\mu^2)\Delta_+(x;\mu^2),\\
 \Delta_+(x;\mu^2)&=\int\frac{\dd^3q}{(2\pi)^3\,2\sqrt{\bm q^2+\mu^2}}
 e^{-\ii q\cdot x},\qquad q^0=\sqrt{\bm q^2+\mu^2}.
\end{align}
Time ordering replaces each positive-frequency function by a Feynman
propagator of the same mass. Thus the K\"all\'en--Lehmann representation is
\begin{equation}
 \boxed{\quad \widetilde G_2(p)=\int_0^\infty\dd\mu^2\,
 \rho(\mu^2)\frac{\ii}{p^2-\mu^2+\ii0},\qquad \rho\geq0.\quad}
 \label{eq:KL}
\end{equation}
Here \(G_2(x)=\intp e^{-\ii p\cdot x}\widetilde G_2(p)\).
If ultraviolet subtractions are needed, local contact terms can be added
to the momentum-space dispersion representation; they do not change its
isolated particle poles or spectral discontinuities.

\begin{keyidea}[A weighted spectral measure]
\(\rho\) measures spectral weight seen by \(\phi\), not merely the number
of states at a given invariant mass. A state with zero field overlap is
absent from this spectral function even if it exists in the theory.
Different operators can therefore reveal different thresholds and poles.
\end{keyidea}

\module{Reading particles and thresholds from the spectrum}
Insert just the one-particle states of \eqref{eq:oneparticle} into
\eqref{eq:rhodef}. Their contribution is
\begin{equation}
 \rho(\mu^2)=\Zp\delta(\mu^2-M^2)+\rho_{\rm cont}(\mu^2),\qquad
 \widetilde G_2(p)=\frac{\ii\Zp}{p^2-M^2+\ii0}
 +\hbox{terms regular near }M^2.
 \label{eq:pole}
\end{equation}
This is the pole-residue statement used in Lecture 22. With the convention
\(\widetilde G_2=\ii/[p^2-m^2-\Sigma_R(p^2)+\ii0]\), a simple stable pole gives
\begin{equation}
 M^2-m^2-\Sigma_R(M^2)=0,\qquad
 \Zp=\frac{1}{1-\Sigma_R'(M^2)}.
\end{equation}
Additional stable particles give additional delta functions if their
quantum numbers and field overlaps permit them.

Thresholds depend on those quantum numbers. In weakly coupled massive
\(\lambda\phi^4\) theory with unbroken \(\phi\mapsto-\phi\), the vacuum is
even and \(\phi\) is odd. Its overlap with a two-\(\phi\) state vanishes.
Assuming no additional lighter states or bound states, the first continuum
in the \(\phi\) two-point function begins at \(\mu^2=9M^2\). It comes from
three particles and first contributes at order \(\lambda^2\).
The even operator \(\phi^2-\langle\phi^2\rangle\), by contrast, can see a
two-particle continuum starting at \(4M^2\).

\begin{center}
\begin{tikzpicture}[x=0.62cm,y=0.65cm,>=stealth]
 \draw[->] (0,0)--(10.7,0) node[right] {$\mu^2$};
 \draw[->] (0,0)--(0,2.55) node[above] {$\rho_\phi$};
 \draw[->,very thick,qftblue] (1,0)--(1,2.1);
 \node[above] at (2.2,2.1) {$\Zp\delta(\mu^2-M^2)$};
 \draw[qftgreen,thick] (9,0) .. controls (9.4,0.05) and (10,0.65) .. (10.5,1.0);
 \node[below] at (1,0) {$M^2$};\node[below] at (9,0) {$9M^2$};
 \node[above] at (8.1,1.35) {odd-sector continuum};
\end{tikzpicture}
\end{center}
\smallskip
The two-particle cut of the \emph{four-point} bubble in Lecture 22 still
starts at \(s=4M^2\). It probes a different channel; there is no conflict.

\begin{remarkbox}[Residue is normalization dependent]
For \(\phi'=c\phi\), with real \(c\ne0\), \(\rho'=c^2\rho\) and
\(Z'_{\rm pole}=c^2\Zp\), while particle masses and scattering probabilities
are unchanged. Positivity implies \(\Zp>0\) for a nonzero overlap.
An upper bound \(\Zp\leq1\) needs an additional canonical spectral sum rule;
it is not a universal statement about an arbitrarily normalized field.
\end{remarkbox}

\module{Incoming and outgoing states: a comparison of dynamics}
Prepare smooth packets with distinct velocities. Far in the past their
centers are widely separated and move toward the collision region; far
in the future the outgoing packets separate again. The interaction
between different packets becomes negligible in these limits. Each
particle remains the same fully interacting, dressed stable excitation.

Denote the maps from the asymptotic particle space to the physical space by
\begin{equation}
 \Omega_{\rm in},\Omega_{\rm out}:\cH_{\rm as}\longrightarrow\cH,
 \qquad
 \ket{\alpha,\mathrm{in/out}}=\Omega_{\rm in/out}\ket{\alpha}_{\rm as}.
\end{equation}
A useful comparison-dynamics notation is
\begin{equation}
 \Omega_{\rm in/out}=\mathop{\mathrm{s\!-!lim}}_{t\to\mp\infty}
 e^{\ii Ht}J e^{-\ii H_{\rm as}t}.
 \label{eq:waveops}
\end{equation}
Here \(J\) identifies separated channel packets with physical states, and
``s-lim'' means convergence in norm on the packet states in its domain.
Equivalently, for such a packet \(f\),
\begin{equation}
 \lim_{t\to\mp\infty}
 \big\|e^{-\ii Ht}\Omega_{\rm in/out}f
       -J e^{-\ii H_{\rm as}t}f\big\|=0.
 \label{eq:comparison}
\end{equation}
It is the \emph{difference} that tends to zero. Neither evolving vector
usually has a limit separately. Nor is \(J\) an identification of the
interacting local field with a bare free field at every time.

In local massive QFT a precise implementation uses mass-shell-filtered
local operators, smeared with packets of disjoint velocity support.
Haag--Ruelle theory constructs the incoming and outgoing limits and their
Fock inner products under appropriate spectral and locality assumptions.
Equation \eqref{eq:waveops} records the comparison principle; an arbitrary
choice of \(J\) is not a proof that these limits exist. We assume the
appropriate construction for the stable channels under study.

\begin{keyidea}[Past and future labels]
Incoming means \(t\to-\infty\); outgoing means \(t\to+\infty\).
The maps obey \(\Omega_{\rm in/out}^\dagger\Omega_{\rm in/out}=1_{\rm as}\)
and intertwine translations, in particular
\(H\Omega_{\rm in/out}=\Omega_{\rm in/out}H_{\rm as}\).
Thus an exact generalized in-state has energy \(\sum_j E_{p_j}\), even
though the state has complicated local correlations during the collision.
\end{keyidea}
Particle separation is not the operation of physically switching the
coupling off. Adiabatic switching is a regulator used in a perturbative
construction; it must not replace the definition of asymptotic particles.

\module{Defining the S-matrix and identifying its Hilbert space}
The amplitude for an incoming configuration \(\alpha\) to be detected as
an outgoing configuration \(\beta\) is
\begin{equation}
 \boxed{\quad S_{\beta\alpha}=\langle\beta,\mathrm{out}|\alpha,\mathrm{in}\rangle,
 \qquad S=\Omega_{\rm out}^\dagger\Omega_{\rm in}.\quad}
 \label{eq:S}
\end{equation}
Both states in the inner product belong to the interacting theory at a
common reference time. The operator \(S\) acts on \(\cH_{\rm as}\): it maps
incoming particle coefficients to outgoing coefficients. If \(f\) and
\(g\) are normalized packet configurations, \(|\langle g|S|f\rangle|^2\)
is an ordinary transition probability for that outgoing packet projection.

Let the ranges of the two wave maps agree and call their common range
\(\cH_{\rm scatt}\), with projector \(P_{\rm scatt}\). Then
\begin{align}
 \Omega_{\rm out}\Omega_{\rm out}^\dagger
 &=\Omega_{\rm in}\Omega_{\rm in}^\dagger=P_{\rm scatt},\\
 S^\dagger S
 &=\Omega_{\rm in}^\dagger P_{\rm scatt}\Omega_{\rm in}=1_{\rm as},
 \qquad SS^\dagger=1_{\rm as}.
 \label{eq:unitarity}
\end{align}
Thus unitarity follows from isometry and equality of the incoming and
outgoing scattering ranges. The stronger assertion
\(\cH_{\rm scatt}=\cH\) is \emph{asymptotic completeness}; it is not a
consequence of writing an overlap or of perturbative diagram rules.

A stable bound state must be included as a species whenever it belongs
to an accessible channel. Omitting it, or omitting particle-production
channels, discards possible outcomes. For an elastic-channel projector
\(P\) on \(\cH_{\rm as}\), set \(S_{\rm el}=PSP\). Then
\begin{equation}
 S_{\rm el}^\dagger S_{\rm el}
 =P-PS^\dagger(1-P)SP.
\end{equation}
The second term is positive. The elastic block need not be unitary even
when the complete S-matrix is unitary: probability can leave that block.

\begin{remarkbox}[Why raw time evolution is not the S-matrix]
Even without interactions, \(e^{-\ii H_{\rm as}(t_f-t_i)}\) contains free
propagation phases. Sending its endpoints to infinity does not give the
identity as an ordinary operator limit. The wave operators remove those
comparison phases. In a genuinely free theory the two wave maps coincide
and \(S=1\), as a scattering experiment requires.
\end{remarkbox}

\module{Symmetry, connected amplitudes, and perturbative calculations}
The wave maps intertwine the Poincar\'e representations. Consequently
\(S\) commutes with the asymptotic total four-momentum. For momentum
configurations with totals \(P_i,P_f\), translation invariance gives
\begin{equation}
 \langle f|S|i\rangle=e^{\ii(P_f-P_i)\cdot a}\langle f|S|i\rangle
 \quad\hbox{for every }a.
\end{equation}
An interacting connected amplitude therefore contains a momentum-conserving
four-dimensional delta function. We define
\begin{equation}
 S=1+\ii T,\qquad
 \langle f|\ii T|i\rangle_c
 =\ii(2\pi)^4\delta^{(4)}(P_f-P_i)\cM_{fi}.
 \label{eq:M}
\end{equation}
The subscript \(c\) excludes spectator overlaps and products of independent
collisions. For an ordinary connected \(2\to2\) process it is usually left
implicit. The identity term is essential: it describes unchanged particles.

For scalar scattering \(\cM\) is Lorentz invariant and depends on invariant
kinematic variables. For \(p_1+p_2\to p_3+p_4\),
\begin{equation}
 s=(p_1+p_2)^2,\quad t=(p_1-p_3)^2,\quad u=(p_1-p_4)^2,
 \qquad s+t+u=\sum_{j=1}^4m_j^2.
\end{equation}
With spin labels, amplitudes transform covariantly with the corresponding
little-group rotations; a single spin component need not be an invariant
scalar. Summed or appropriately averaged probabilities have the expected
Lorentz-invariant description.

The definitions above do not assume small coupling. Perturbation theory
is a method for calculating them. With a regulator and an interaction-picture
split, time evolution is represented by
\begin{equation}
 U_I(t_f,t_i)=\Tord\exp\left[-\ii\int_{t_i}^{t_f}\dd t\,H_I(t)\right].
\end{equation}
Vacuum normalization, physical external poles, and their residues are
needed to extract the physical amplitude from the resulting correlators.
The path-integral ratio in Lecture 21 removes vacuum bubbles; it does not
remove disconnected external pairings. Lecture 25 supplies the precise
LSZ reduction instead of treating the Dyson exponential as a definition
that bypasses these issues.

\begin{keyidea}[Three distinct objects]
A correlator may be off shell and contain external propagators.
An amputated connected kernel has those propagators removed.
The physical amplitude additionally uses on-shell external momenta and
the correct pole-residue normalization. At tree level in our scalar
quartic theory, these steps give \(\ii\cM=-\ii\lambda\), hence
\(\cM=-\lambda\).
\end{keyidea}

\module{Unitarity as a sum over physical outcomes}
Insert \(S=1+\ii T\) into \(S^\dagger S=1\). The result is
\begin{equation}
 -\ii(T-T^\dagger)=T^\dagger T.
 \label{eq:Tunitarity}
\end{equation}
For any normalized incoming packet, the diagonal matrix element is a sum
of nonnegative transition probabilities. In a plane-wave basis, insert
\eqref{eq:complete}, factor out the common momentum delta function, and
obtain the forward optical theorem
\begin{equation}
 2\operatorname{Im}\cM_{ii}
 =\sum_X\frac{1}{\cS_X}\sum_{\rm spins}
 \int\dd\Phi_X(P_i)\,|\cM_{Xi}|^2,
 \label{eq:optical}
\end{equation}
\begin{equation}
 \dd\Phi_n(P)=(2\pi)^4\delta^{(4)}\!\left(P-\sum_{j=1}^np_j\right)
 \prod_{j=1}^n\frac{\dd^3p_j}{(2\pi)^3\,2E_j}.
 \label{eq:phase}
\end{equation}
For two incoming stable particles in a massive theory, the amplitudes on
the right describe all open scattering channels. Initial spin labels are
fixed unless we explicitly average. A fully general multiparticle matrix
identity retains spectator and disconnected terms as well.

Why does a real tree amplitude not violate \eqref{eq:optical}? If
\(\cM^{(1)}\) is first order in a coupling, the right side starts at second
order. At that order,
\begin{equation}
 2\operatorname{Im}\cM^{(2)}_{ii}
 =\sum_X\frac1{\cS_X}\int\dd\Phi_X\,
 |\cM^{(1)}_{Xi}|^2.
\end{equation}
A loop imaginary part supplies the probability required by the tree
transition. This explains the two-particle cut encountered in Lecture 22
without redoing its loop integral. Lecture 24 turns the same identity into
a relation involving the measurable total cross section.

\begin{remarkbox}[Where the construction needs modification]
Stable massive particles with isolated poles are the clean setting of
these lectures. Long-range forces and soft massless radiation can obstruct
the ordinary Fock-space limits; charged states in QED require additional
infrared care. An unstable particle has no exact real-axis one-particle
pole suitable for an external LSZ state. Neither complication means that
experiments lack predictions; it means the asymptotic description must be
adapted to the observable.
\end{remarkbox}
\textbf{Checks before moving on.} Explain why \(\phi\ket{\Omega}\) need
not be a one-particle state, why both \(\Omega\) maps are needed, and why
unitarity requires a sum over all accessible channels. These are the
physical inputs behind the rate formulas in the next lecture.
\end{document}
